% answers_short.tex —— 1990 年上海卷数学 · 简短答案速查
% 由 tools/make_short_answers.py 自 answers.tex 的 \gksolhead/\gkanswer 提取，
% 并与 1990/1990_数学真题_上海_答案要点.md 逐题比对后生成，勿手改（改源 MD 后重跑）。

\begin{gknotice}
\textbf{使用说明}：本节只给最终答案便于快速核对，完整推导见解析版。答案为编写者独立推导，\textbf{非官方标准答案}。
\end{gknotice}

\begin{gknotice}
选择题答案中的 U 表示待核，未给出猜测答案，详情见解析版。
\end{gknotice}

\gkpart{填空题}
\begin{itemize}[leftmargin=2.6em]
  \item[\textbf{1.}] $[-4,+\infty)\setminus\{-2\}$
  \item[\textbf{2.}] $y=\sin x$，$x\in[-\frac\pi2,\frac\pi2]$
  \item[\textbf{3.}] $2x+y-4=0$
  \item[\textbf{4.}] $2\pi\,\mathrm{cm}^3$
  \item[\textbf{5.}] $\frac{2}{\sqrt5}$
  \item[\textbf{6.}] $0$
  \item[\textbf{7.}] $al$
  \item[\textbf{8.}] $a=-2$
  \item[\textbf{9.}] $m=-\frac43$
  \item[\textbf{10.}] $60$
\end{itemize}

\gkpart{选择题}
\begin{center}\begin{tabular}{ccccccccccc}
\toprule
题号 & 11 & 12 & 13 & 14 & 15 & 16 & 17 & 18 & 19 & 20 \\
\midrule
答案 & C & B & A & A & C & B & B & U & B & D \\
\bottomrule
\end{tabular}\end{center}

\gkpart{解答题}
\begin{itemize}[leftmargin=2.6em]
  \item[\textbf{21.}] $y=9\sqrt5$
  \item[\textbf{22.}] $x=2\pi-\arccos\frac13$
  \item[\textbf{23.}] 在 $P$ 位于直线 $x=2$ 的条件下，轨迹为 $4(x-\frac12)^2+2y^2=1$（不含 $(1,0)$），焦点为 $(\frac12,\pm\frac12)$
  \item[\textbf{24.}] $2$
  \item[\textbf{25.}] $a\in\{-1\}\cup[1,3]$
  \item[\textbf{26.}] (1) $\frac{\sqrt6}{4}$；(2) $\arccos\sqrt{\frac35}$
  \item[\textbf{27.}] $\pi-\arctan\frac{5}{2\sqrt6}$
\end{itemize}
